% =============================================================================
% 038 / 068
% Status: raw
% =============================================================================
% Notes:
%
% Source question:
% 請問小字與吳語拼音及某吳語方言音系的關系是
% 實際音系→吳語拼音→小字
% 還是
% 吳語拼音←實際音系→小字？
% (盡管說吳拼是音位拼音，但由於不同地區音系差異，吳拼方案總會有區別，尤其是與內部一致性高的北吳所不同的南吳。)
% =============================================================================

\chapter[請問小字與吳語拼音及某吳語方言音系的關系是 實際音系→吳語拼音→小字 還是 吳語拼音←實際音系→小字？ (盡管說吳拼是音位拼音，但由於不同地區音系差異，吳拼方案總會有區別，尤其是與內部一致性高的北吳所不同的南吳。)]{請問小字與吳語拼音及某吳語方言音系的關系是 \\[0.35em] 實際音系→吳語拼音→小字 \\[0.35em] 還是 \\[0.35em] 吳語拼音←實際音系→小字？ \\[0.35em] (盡管說吳拼是音位拼音，但由於不同地區音系差異，吳拼方案總會有區別，尤其是與內部一致性高的北吳所不同的南吳。)}
\qaanswer{
漢字以音位拼音為主，各地可以先按自己的發音習慣來，這樣7成以上就是一個體系的，音值不同沒關係。等有了一定書面語創作的量就看哪個地區的創作者多了。露出度高的地區的發音音值習慣就有競爭力。這個得看自然演化。如果書面語都出不來，吳語完球是無法避免的。
}

